\subsection*{J 短期冲击模型补充检验}

正文“模型建立与求解”章只保留校准、基准对比和鲁棒性摘要。附录补充给出短期模型的机制贡献、参数剖面、事件窗口、滞后平移、残差分布和局部扰动稳健性图。

\appendixfigunit{0.39\textheight}{机制贡献分解}{短期价格平台不是由单一变量决定，而是物理缺口、风险溢价、政策缓冲和预期修复共同作用。该图用于说明各机制在不同日期上的方向和相对强度。}{0.96}{0.29\textheight}{短期模型机制贡献.png}{附图 1 机制贡献分解}

\appendixfigunit{0.39\textheight}{参数剖面检验}{参数剖面图观察关键参数在最优点附近扰动时误差函数是否平滑变化。若曲线只在某一点突然变好，就可能存在单点调参风险。}{0.96}{0.30\textheight}{短期模型参数剖面图.png}{附图 2 参数一维剖面}

附图 1 和附图 2 分别回答“价格平台由哪些机制推动”和“参数是否只在单点有效”。前者给出机制方向，后者观察目标函数形态，二者共同支撑短期模型的解释性和稳定性。

\appendixfigunit{0.39\textheight}{事件窗口检验}{短期冲击窗口内部并不均匀，初始冲击、高位平台、中期回落和后期二次抬升的拟合难度不同。该图把整体误差拆到事件阶段，避免只用总 RMSE 掩盖局部表现。}{0.84}{0.23\textheight}{短期模型事件窗口检验.png}{附图 3 事件窗口检验}

\appendixfigunit{0.36\textheight}{滞后平移检验}{时间序列模型容易出现“把昨天当今天”的滞后拟合。该检验将预测曲线进行平移比较；原始位置误差最低，说明模型不是简单复制滞后价格路径。}{0.78}{0.21\textheight}{短期模型滞后平移检验.png}{附图 4 滞后平移检验}

\appendixfigunit{0.39\textheight}{残差分段分布}{残差分布用于判断误差集中在哪些阶段。中期平台阶段残差更分散，说明平台震荡和情绪再定价仍是短期模型的主要误差来源。}{0.84}{0.23\textheight}{短期残差分段分布图.png}{附图 5 残差分段分布}

\appendixfigunit{0.52\textheight}{局部扰动稳健性}{稳健性带展示当前参数邻域内的预测路径范围。若扰动后路径仍围绕真实价格平台展开，说明结论不依赖某一个极窄的参数组合。}{1.00}{0.43\textheight}{短期模型稳健性带.png}{附图 6 局部扰动稳健性}

以上六组补充检验依次对应机制来源、参数形态、阶段误差、滞后排查、残差分布和局部扰动。它们共同把正文的短期拟合结论拆成可复核证据链，说明模型优势来自机制组合，而不是单点调参或滞后复制。

\clearpage
